\documentclass[12pt,a4paper,oneside]{article}
\usepackage[brazil]{babel}
\usepackage[utf8]{inputenc}
\usepackage{olymp}


\newlength{\exmpwidouf}
\exmpwidinf=10cm
\exmpwidouf=6cm


\setcounter{problem}{0}

\begin{document}

\begin{problem}{Idade da vovó Mônica}{idade}{2}

Vovó Mônica é avó de vários netos. Ela notou que, neste ano, a soma das idades dos seus netos é igual à idade dela. Neste problema, dada a idade de vovó Mônica e as idades de seus netos, exceto um, seu programa deve computar e imprimir a idade do neto mais velho.

Pos exemplo, se sabemos que vovó Mônica tem 52 anos e 3 netos, e as idades conhecidas de dois de seus netos são 14 e 18 anos, então a idade do outro neto, que não era conhecida, tem que ser 20 anos, pois a soma das três idades tem que ser 52. Portanto, a idade do neto mais velho é 20. Em mais um exemplo, se vovó Mônica tem 47 anos e 4 netos, e as idades de três deles são 12, 10 e 14 anos, então o outro neto tem que ter 11 anos e, portanto, a idade do neto mais velho é 14.

%\Notes
%
%Observações, se tiver.

\InputFile

A primeira linha da entrada contém um inteiro $M$ representando a idade de vovó Mônica. A segunda linha contém um inteiro $N$ representando o número de netos de vovó Mônica. a terceira linha contém $N-1$ inteiros, representando as idades conhecidas de seus netos.

Restrições: $40 \leq M \leq 110$, $1 < N \leq 10$.

\OutputFile

Seu programa deve imprimir uma linha, contendo um número inteiro, representando a idade do neto mais velho de vovó Mônica.

% \Notes
% Preste atenção aos tipos das variáveis, pois $F(90) \approx 2^{61}$.

\Example

\begin{example}
\exmp{
52
3
14 18
}{
20
}%
\end{example}

\begin{example}
\exmp{
47
4
12 10 14
}{
14
}%
\end{example}

\begin{example}
\exmp{
40
4
10 10 10 
}{
10
}%
\end{example}

\begin{example}
\exmp{
40
2
19
}{
21
}%
\end{example}

\textit{Problema baseado na questão ``A idade de dona Mônica'' da OBI2019 fase local.}

\end{problem}

\end{document}
